\chapter{Analisi dei Requisiti}
\label{ch:requisiti}

In questo capitolo vengono descritti i requisiti funzionali e non funzionali che il sistema ChillChain è chiamato a soddisfare. L'analisi è stata condotta considerando i tre profili utente previsti dalla piattaforma --- amministratore, cliente e tecnico.

\section{Requisiti Funzionali}

I requisiti funzionali definiscono le capacità operative che la piattaforma deve offrire ai propri utenti e descrivono il comportamento atteso del sistema in risposta a specifiche azioni o eventi.

\subsection{Gestione degli utenti e autenticazione}

Il sistema deve consentire l'autenticazione mediante credenziali protette. A ogni utente autenticato viene assegnato un ruolo tra amministratore e tecnico. Il cliente non ha bisogno di autenticazione.

\subsection{Gestione della flotta}

L'amministratore deve poter registrare nuovi veicoli nella piattaforma, specificando per ciascuno le dimensioni del vano di carico (larghezza, altezza, lunghezza), la portata massima in peso, il costo per chilometro e l'indicazione sulla presenza o meno dell'impianto di refrigerazione. Il sistema deve consentire la visualizzazione dell'intera flotta e la rimozione dei veicoli non più operativi.

\subsection{Gestione dei viaggi e delle spedizioni}

L'amministratore deve poter visualizzare l'elenco dei viaggi programmati e consultare le spedizioni associate a ciascun viaggio.

\subsection{Ricerca e prenotazione delle spedizioni}

Il cliente deve poter ricercare i viaggi disponibili inserendo l'indirizzo di partenza, l'indirizzo di arrivo, le caratteristiche del proprio pacco --- dimensioni, peso, necessità di refrigerazione --- e la data entro cui la consegna deve essere completata. 
Il sistema deve restituire due categorie di risultati: nuovi viaggi, generati a partire dai veicoli della flotta ancora privi di viaggi assegnati con percorso calcolato ad hoc, e viaggi condivisi, corrispondenti a viaggi già programmati che dispongano di capacità residua sufficiente e di compatibilità geografica con il percorso richiesto. Per i viaggi condivisi, il prezzo deve essere ricalcolato applicando un meccanismo di divisione basato sul numero di spedizioni già presenti, in modo da incentivare la condivisione del veicolo. Il cliente seleziona il viaggio desiderato tra quelli proposti. Alla conferma della selezione, il sistema deve: nel caso di un nuovo viaggio, crearlo automaticamente con il percorso calcolato; nel caso di un viaggio condiviso, ricalcolare dinamicamente il percorso introducendo le fermate del nuovo cliente come tappe intermedie.

\subsection{Monitoraggio telemetrico in tempo reale}

Il sistema deve acquisire i dati telemetrici provenienti dai sensori installati a bordo dei veicoli refrigerati. Il tecnico deve poter visualizzare i dati telemetrici attraverso una dashboard che presenti l'andamento temporale dei parametri mediante grafici, con la possibilità di filtrare per veicolo e per tipologia di dato telemetrico.

\subsection{Rilevazione delle anomalie}

Un servizio dedicato deve analizzare in maniera continuativa i flussi telemetrici provenienti dai veicoli e identificare automaticamente condizioni anomale. Quando un'anomalia viene confermata, il sistema deve generare una notifica persistente, associata al veicolo interessato.
\subsection{Gestione delle notifiche}

Il tecnico deve poter consultare l'elenco delle notifiche ricevute, ciascuna contenente il nome del veicolo, il timestamp dell'evento, l'errore di ricostruzione e un messaggio descrittivo dell'anomalia rilevata. 
\section{Requisiti Non Funzionali}

I requisiti non funzionali specificano le qualità architetturali e operative che il sistema deve garantire trasversalmente rispetto alle singole funzionalità.

\subsection{Usabilità}

L'interfaccia utente deve essere progettata con particolare attenzione alla chiarezza e all'immediatezza d'uso. Le informazioni devono essere presentate in modo differenziato in base al ruolo dell'utente.

\subsection{Sicurezza}

Il sistema deve proteggere le comunicazioni tra client e server mediante autenticazione e assegnazione dei ruoli agli utenti autenticati. Le password degli utenti devono essere memorizzate in forma cifrata. 

\subsection{Modularità e manutenibilità}

L'architettura del sistema deve seguire il paradigma a microservizi, con una separazione netta tra le componenti applicative. Ogni componente deve poter essere sviluppato, testato e aggiornato indipendentemente dagli altri.

\subsection{Portabilità e riproducibilità del deployment}

L'intera infrastruttura deve essere containerizzata tramite Docker e orchestrata mediante Docker Compose, in modo da garantire che il sistema possa essere avviato in maniera identica su qualsiasi macchina che disponga del runtime Docker.
